\section{费米液体的平衡性质（Gross 视角）}

Gross 第 31 章强调：一旦给定 $\{F_\ell\}$，平衡热力学由局域平衡分布的形变完全确定。宏观公式已见于 \eqref{eq:mstar}--\eqref{eq:Kchi_FL}；此处补形变语言与推导要点。

\subsection{准粒子分布的形变}

平衡：$n_{\mathbf{p}}^0=\theta(-\xi_{\mathbf{p}})$（$T=0$）或费米函数。弱扰动下局域平衡要求
\begin{equation}
  \delta n_{\mathbf{p}}
  =-\frac{\partial n^0}{\partial\varepsilon_{\mathbf{p}}}
  \bigl(
  \delta\varepsilon_{\mathbf{p}}-\delta\mu-\mathbf{p}\cdot\mathbf{u}-\cdots
  \bigr),
\end{equation}
其中 $\delta\varepsilon_{\mathbf{p}}=\sum_{\mathbf{p}'}f_{\mathbf{p}\mathbf{p}'}\delta n_{\mathbf{p}'}$。把 $\delta n$ 在费米面按球谐展开，各 $\ell$ 通道由 $F_\ell^{s,a}$ 闭合。

\subsection{从形变到响应}

\begin{itemize}
  \item 化学势/密度扰动 $\Rightarrow$ $\ell=0$ 对称通道 $\Rightarrow$ 压缩率 $\kappa/\kappa_0=(m^*/m)/(1+F_0^s)$；
  \item 自旋依赖化学势 $\Rightarrow$ 反对称 $F_0^a$ $\Rightarrow$ $\chi/\chi_0=(m^*/m)/(1+F_0^a)$；
  \item 流形变（伽利略）$\Rightarrow$ $F_1^s$ $\Rightarrow$ $m^*/m=1+F_1^s/3$。
\end{itemize}
比热 $C_V=\frac{m^*}{m}C_V^{(0)}$。Pomeranchuk 稳定性：$1+F_\ell^{s,a}/(2\ell+1)>0$。

\begin{note}
这些强度量对应密度/自旋响应在 $q\to0$、$\omega/(qv_F)\to0$ 的静态极限，与第 5 章线性响应及微观 Landau 极限一致。
\end{note}
